\documentclass[../General.tex]{subfiles}
\begin{document}

\onehalfspacing
\fancypagestyle{plain}{
\fancyhf{}
\rhead{\thepage}
\renewcommand{\headrulewidth}{0pt}
\renewcommand{\footrulewidth}{0pt}
}

%----------------------------------------------------------------------------------------
%	Capítulo 6 - DISEÑO ELECTRÓNICO
%----------------------------------------------------------------------------------------

\chapter{DISEÑO ELECTRÓNICO}

%----------------------------------------------------------------------------------------
%   INDÍCE DEL CAPÍTULO:
%----------------------------------------------------------------------------------------
%   - Texto introductorio -> Presentación de arquitectura electrónica (diagrama de bloques)
%\subfile{cap6/601_arquitectura_del_sistema}
%
%   - Subsistema de sensores e instrumentación meteorológica
%\subfile{cap6/602_sensores_e_instrumentacion}
%
%       - Análisis y dimensionamiento de componentes.
%       - Selección de componentes: Proceso de selección de componentes específicos, tablas comparativa entre alternativas, explicación de criterios.
%       - Bloque de diagrama esquemático correspondiente y disposición inmediata de subsistema. [1]
%
%   - Subsistema de actuadores y controladores de potencia
%\subfile{cap6/603_actuadores_y_potencia}
%       - Análisis y dimensionamiento de componentes.
%       - Selección de componentes: ""
%       - Bloque de diagrama esquemático correspondiente y disposición inmediata de subsistema.
%
%   - Subsistema de comunicación
%\subfile{cap6/604_comunicacion}
%       - Arquitectura de telemetría de largo alcance.
%       - Análisis y dimensionamiento de componentes.
%       - Selección de componentes: ""
%       - Bloque de diagrama esquemático correspondiente y disposición inmediata de subsistema.
%
%   - Subsistema de procesamiento y control
%\subfile{cap6/605_procesamiento_y_control}
%       - Arquitectura de control: Controladores usados y comunicación entre controladores.
%       - Análisis y dimensionamiento de componentes.
%       - Selección de componentes: ""
%       - Bloque de diagrama esquemático correspondiente y disposición inmediata de subsistema.
%
%   - Subsistema de energía
%\subfile{cap6/606_energia}
%       - Balance energético del sistema
%       - Análisis y dimensionamiento de componentes.
%       - Selección de componentes: ""
%       - Bloque de diagrama esquemático correspondiente y disposición inmediata de subsistema.
%       - Diseño de PCB: Parámetros de placa, disipación térmica, dimensiones, imagen del CAD de la PCB.
%   - Integración y distribución de subsistemas electrónicos
%\subfile{cap6/607_integracion_y_distribucion}
%
%
% ANEXOS: 
% - Cálculos y explicación de componentes adyacentes a cada subsistema (resistencias, capacitores, inductores alrededor de los ICs).
% - Diagrama esquemático completo.
% - BOM de componentes electrónicos.
% NOTAS:
% [1]: La disposición inmediata es la representación del espacio físico (incluyendo soportes y los mismos módulos) AISLADO del subsistema, no incluye sus conexiones a otros subsistemas, si hay otros subsistemas cerca de este o en el mismo compartimento, se invisibilizan (piensa en la vista del sistema nervioso o el sistema respiratorio de un cuerpo).


\end{document}